\documentclass[11pt]{amsart}
\usepackage{amsmath,amssymb,amsthm}
\usepackage[margin=1.1in]{geometry}
\usepackage[hidelinks]{hyperref}
\usepackage{etoolbox}
\usepackage{xcolor}

% amsart prints addresses after the references; put the affiliation under the author on the title page.
\makeatletter
\patchcmd{\@setauthors}{\MakeUppercase{\authors}}%
  {\MakeUppercase{\authors}\par\vspace{3pt}{\normalfont\small AIM Intelligence}}%
  {}{\PackageError{note}{could not place the affiliation}{}}
\makeatother

% Unicode characters of the Lean listing in Section 6.
\DeclareUnicodeCharacter{2115}{\ensuremath{\mathbb{N}}}
\DeclareUnicodeCharacter{00D7}{\ensuremath{\times}}
\DeclareUnicodeCharacter{2264}{\ensuremath{\le}}
\DeclareUnicodeCharacter{2203}{\ensuremath{\exists}}
\DeclareUnicodeCharacter{2227}{\ensuremath{\wedge}}
\DeclareUnicodeCharacter{2200}{\ensuremath{\forall}}
\DeclareUnicodeCharacter{2192}{\ensuremath{\to}}

\newtheorem{theorem}{Theorem}
\newtheorem{lemma}[theorem]{Lemma}
\newtheorem{corollary}[theorem]{Corollary}
\theoremstyle{remark}
\newtheorem{remark}[theorem]{Remark}

\newcommand{\F}{\mathbb{F}}
\newcommand{\Q}{\mathbb{Q}}
\newcommand{\rk}{\operatorname{rank}}
\newcommand{\R}{\mathcal{R}}
\newcommand{\col}{\operatorname{col}}

\title{Lower bounds for Naslund's crossing problem}
\author{Arth Singh}
\thanks{{\color{red}We share this mathematics as a contribution to the community. We do not intend to
affect anyone's ongoing work in this direction.}}
\date{}

\begin{document}

\begin{abstract}
Let $T_n(x_1,x_2,x_3,x_4)=[x_1=x_4]\,[x_2=x_3]$ on $[n]^4$ over a field $\F$. Naslund asked whether $T_n$
is a sum of $r$ tensors of rank one along the partition $\{12\mid34\}$ and $s$ tensors of rank one along
$\{13\mid24\}$ only if $r+s\ge n^{1+\varepsilon}$. In matrix form, the least such $r+s$ is the least value
of $\rk B+\rk\R(C)$ over all splittings $S_n=B+C$ of the swap matrix of $\F^n\otimes\F^n$, where $\R$ is the
realignment. The only lower bound we found in the literature is the trivial bound $n$. We prove
$r+s\ge 4n/3$ for $n\ge3$, $r+s\ge15n/11$ for $n\ge5$, $r+s\ge67n/49$ for $n\ge7$ and $r+s\ge n+3$ for
$n\ge4$, over every field. The least value is $4$ at $n=3$ and
$7$ at $n=4$; at $n=4$, limits of splittings reach $6$ in characteristic $0$ and in all but seven positive
characteristics. The proof of the bound $4n/3$ takes one page. The
lower bounds and the value at $n=4$ are formalized in Lean~4 with Mathlib. The theorems and their formal proofs were found by AI agents
in a multi-agent research harness.
\end{abstract}

\maketitle

\section{Introduction}

Let $\F$ be a field and $n\ge1$. Write $[P]$ for $1$ if the statement $P$ is true and $0$ otherwise. The
\emph{crossing} is the tensor
\[
  T_n(x,y,u,v)=[x=v]\,[y=u],\qquad x,y,u,v\in[n].
\]
A tensor $f\colon[n]^4\to\F$ has rank one along $\{12\mid34\}$ if $f(x,y,u,v)=a(x,y)\,b(u,v)$ for some
functions $a,b$, and rank one along $\{13\mid24\}$ if $f(x,y,u,v)=a(x,u)\,b(y,v)$. Let $P_\F(n)$ be the least
$r+s$ such that $T_n$ is a sum of $r$ tensors of the first kind and $s$ tensors of the second kind. This is
the rank of $T_n$ with respect to the family of partitions $\{\{12\mid34\},\{13\mid24\}\}$ in Naslund's
framework of partition ranks~\cite{Nas20}.

Naslund asked~\cite[Problem~12]{Nas20} whether $P_\F(n)\ge n^{1+\varepsilon}$ for some $\varepsilon>0$, and
noted that no explicit superlinear lower bound is known for any tensor of this kind. Karam~\cite[Question~1.2]
{Kar23} restates the question and records the lower bound $n$: if two of the four coordinates are fixed, an
$n\times n$ identity matrix remains, and it is a sum of at most $r+s$ matrices of rank one. A positive answer
would give new lower bounds for non-commutative circuits~\cite{HWY11,Nas20}. In the same framework, Brand,
Kaski and Wang~\cite{BKW26} use partition ranks to bound multiplicative complexity from below in any constant
degree; their paper does not discuss the crossing tensor.

\emph{Matrix form.} Index the rows of an $n^2\times n^2$ matrix by pairs $(x,y)$ and its columns by pairs
$(u,v)$. The swap matrix $S_n$ has entries $S_n[(x,y),(u,v)]=[u=y]\,[v=x]=T_n(x,y,u,v)$. The
\emph{realignment} of a matrix $C$ is the matrix $\R(C)$ with $\R(C)[(x,u),(y,v)]=C[(x,y),(u,v)]$; its rank
is the operator Schmidt rank of $C$. A sum of $r$ tensors of the first kind is a matrix of rank at most $r$,
and a sum of $s$ tensors of the second kind is a matrix $C$ with $\rk\R(C)\le s$; the converses also hold.
Hence
\[
  P_\F(n)=\min\{\rk B+\rk\R(C) : S_n=B+C\}.
\]
The realignment is an involution and $\R(S_n)=S_n$, so $S_n=\R(C)+\R(B)$ is again a splitting, with the roles
of $\rk B$ and $\rk\R(C)$ exchanged.

\subsection*{Results}

\begin{theorem}\label{thm:43}
Let $\F$ be a field and $n\ge3$. If $S_n=B+C$, then $3\,(\rk B+\rk\R(C))\ge4n$.
\end{theorem}

\begin{theorem}\label{thm:1511}
Let $\F$ be a field and $n\ge5$. If $S_n=B+C$, then $11\,(\rk B+\rk\R(C))\ge15n$.
\end{theorem}

\begin{theorem}\label{thm:6749}
Let $\F$ be a field and $n\ge7$. If $S_n=B+C$, then $49\,(\rk B+\rk\R(C))\ge67n$.
\end{theorem}

\begin{theorem}\label{thm:n3}
Let $\F$ be a field and $n\ge4$. If $S_n=B+C$, then $\rk B+\rk\R(C)\ge n+3$.
\end{theorem}

\begin{corollary}\label{cor:values}
For every field $\F$, $P_\F(n)\ge\max(\lceil 15n/11\rceil,\,n+3)$ for $n\ge5$, and
$P_\F(n)\ge\max(\lceil 67n/49\rceil,\,\lceil 15n/11\rceil,\,n+3)=\lceil 67n/49\rceil$ for $n\ge7$. Moreover $P_\F(3)=4$,
$P_\F(4)=7$ and $8\le P_\F(5)\le11$.
\end{corollary}

To our knowledge, these are the first lower bounds on $P_\F(n)$ above the trivial bound $n$. They are far
from Naslund's question: all four are linear in $n$. Theorem~\ref{thm:43} is attained at $n=3$, and its
proof takes one page. Theorems~\ref{thm:1511} and~\ref{thm:6749} have the better constants
$67/49>15/11>4/3$; their proofs are certificates that a computer search found. Theorem~\ref{thm:n3} is the
best of the four bounds for $n=4$ and $n=5$. For $n\ge6$ it follows from Theorem~\ref{thm:1511}, because
then $\lceil15n/11\rceil\ge n+3$, so its new content is the two cases $n=4$ and $n=5$. The hypotheses on $n$
are needed: $P_\F(2)=2$, so Theorem~\ref{thm:43} fails at $n=2$, and $P_\F(3)=4$, so
Theorems~\ref{thm:1511}, \ref{thm:6749} and~\ref{thm:n3} fail at $n=3$. The inequalities of
Theorems~\ref{thm:1511} and~\ref{thm:6749} also hold for the smaller values $n\ge4$, by
Theorem~\ref{thm:n3}.

\emph{Connection to circuit lower bounds.} Over $\Q$, $P_\Q(n)$ is the width of the non-commutative
polynomial $\sum_{i,j\in[n]}x_iy_jx_iy_j$ \cite[Lemma~3.5]{HWY11}. By \cite[Proposition~4.1]{HWY11} it lies
between $\beta(n)$ and $4\beta(n)$, where $\beta(n)$, written $\beta_\Q(n)$ in~\cite{Hru24}, is the least
number of terms in an identity $(x_1^2+\dots+x_n^2)(y_1^2+\dots+y_n^2)=\sum_{k=1}^{\beta(n)}f_kf_k'$ with
bilinear forms $f_k,f'_k$ over $\Q$; see also~\cite[Appendix~G]{Cla26}. A lower bound
$\beta(n)\ge n^{1+\varepsilon}$ would give exponential lower bounds on non-commutative circuits over $\Q$ for
the permanent~\cite[Theorem~1.8]{HWY11}, and any superlinear lower bound on $\beta(n)$ would give a
superlinear lower bound on non-commutative circuits for an explicit polynomial of degree four
\cite[Theorem~1.7]{HWY11}. For the least number $S_{\mathbb R}(n)$ of squares of real bilinear forms in such
an identity, the bounds $S_{\mathbb R}(n)>n$ unless $n\in\{1,2,4,8\}$ and $S_{\mathbb R}(n)\ge(2-o(1))n$ are
known \cite[Theorems~1.1 and~1.2]{HWY11}. They use that the terms are squares, and they do not apply to
$\beta(n)$ or to $P_\F(n)$. The best upper bounds we know are $P_\Q(n)=O(n^{11/6})$
\cite[Appendix~G]{Cla26} and, through $P_\Q(n)\le4\beta(n)$, $P_\Q(n)=O(n^{1.62})$
by~\cite[Theorem~12]{Hru24}, which holds over every field; its proof is sketched there as a modification
of~\cite[Theorem~1]{Hru24}.

\emph{Ideas of the proofs.} A splitting $S_n=B+C$ with $r=\rk B$ and $s=\rk\R(C)$ gives, for every pair of
indices $y,u\in[n]$, an identity $X_yY_u^{\top}+Z_uW_y^{\top}=[y=u]\,I_n$ between $n\times r$ and $n\times s$
matrices (Lemma~\ref{lem:slices}). For $y=u$ this gives the trivial bound $r+s\ge n$. For $y\ne u$ the
identity forces the column spaces of $X_y$ and $Z_u$ to overlap, by an amount that the ranks of $Y_u$ and
$W_y$ control (Lemma~\ref{lem:excess}). For Theorem~\ref{thm:43} we combine these overlaps around a cycle of
three indices with the submodularity of dimension (Lemma~\ref{lem:cycle}), and we add the resulting
inequalities over the three rotations of the cycle, for the identities and for their transposes.
Theorem~\ref{thm:1511} combines inequalities of the same kinds for five indices; the combination has $189$
terms and was found by a linear-programming search. Theorem~\ref{thm:6749} does the same with seven
indices. For Theorem~\ref{thm:n3}, Theorem~\ref{thm:43} leaves only $r+s=n+2$, and we exclude it by a case analysis on
the pair $(r,s)$.

\emph{Formal verification and provenance.} Theorems~\ref{thm:43}, \ref{thm:1511},
\ref{thm:6749} and~\ref{thm:n3} and the value $P_\F(4)=7$ are proved in Lean~4 with Mathlib (Section~\ref{sec:lean}). The theorems and their Lean proofs
were found by AI agents in one run of a multi-agent harness; Section~\ref{sec:how} says how. The proof of
Theorem~\ref{thm:43} is given in full in Section~\ref{sec:43}. The proofs of Theorems~\ref{thm:1511},
\ref{thm:6749} and~\ref{thm:n3} are only described here, and no person has read them in full: their only complete proofs
are the Lean proofs. For this reason
Section~\ref{sec:lean} prints the formal definitions and statements.

\emph{Organization.} Section~\ref{sec:slices} sets up the identities between slices.
Section~\ref{sec:43} proves Theorem~\ref{thm:43} and describes the proofs of Theorems~\ref{thm:1511}
and~\ref{thm:6749}. Section~\ref{sec:n3} outlines the proof of
Theorem~\ref{thm:n3}, whose formal version is long. Section~\ref{sec:small} proves
Corollary~\ref{cor:values} and makes some remarks, among them a border splitting at $n=4$ with $r+s=6$.

\section{Slices}\label{sec:slices}

Throughout, $S_n=B+C$, $r=\rk B$ and $s=\rk\R(C)$. Choose decompositions into terms of rank one,
\begin{align*}
  B[(x,y),(u,v)]&=\sum_{i=1}^{r}U_i(x,y)\,V_i(u,v),\\
  C[(x,y),(u,v)]=\R(C)[(x,u),(y,v)]&=\sum_{l=1}^{s}G_l(x,u)\,H_l(y,v).
\end{align*}
For $y,u\in[n]$ define the matrices $X_y,Y_u\in\F^{n\times r}$ and $Z_u,W_y\in\F^{n\times s}$ by
\[
  (X_y)_{x,i}=U_i(x,y),\quad (Y_u)_{v,i}=V_i(u,v),\quad (Z_u)_{x,l}=G_l(x,u),\quad (W_y)_{v,l}=H_l(y,v).
\]

\begin{lemma}[slices]\label{lem:slices}
For all $y,u\in[n]$, $X_yY_u^{\top}+Z_uW_y^{\top}=[y=u]\,I_n$.
\end{lemma}

\begin{proof}
The $(x,v)$ entry of the left side is
\[
  \sum_{i}U_i(x,y)V_i(u,v)+\sum_{l}G_l(x,u)H_l(y,v)=B[(x,y),(u,v)]+C[(x,y),(u,v)]=[u=y]\,[v=x].
  \qedhere
\]
\end{proof}

We call a family $(X_y,Y_u,Z_u,W_y)_{y,u\in[n]}$ of $n\times r$ and $n\times s$ matrices with the property of
Lemma~\ref{lem:slices} a \emph{slice system} of widths $(r,s)$. The transpose of the identity for $(u,y)$ is
$Y_yX_u^{\top}+W_uZ_y^{\top}=[y=u]\,I_n$, so $(Y_y,X_u,W_u,Z_y)$ is again a slice system of widths $(r,s)$;
we call it the \emph{transposed} system.

For $y=u$, Lemma~\ref{lem:slices} writes $I_n$ as the product of the $n\times(r+s)$ matrix $[X_y\mid Z_y]$
and the transpose of $[Y_y\mid W_y]$. Hence $r+s\ge n$, which is the trivial bound, and
\begin{equation}\label{eq:span}
  \col X_y+\col Z_y=\F^n\qquad\text{for every }y\in[n],
\end{equation}
where $\col$ denotes the column space. In particular $\rk X_y+\rk Z_y\ge n$, and in the same way
$\rk Y_y+\rk W_y\ge n$.

\section{Proof of Theorem~\ref{thm:43}}\label{sec:43}

Write $\xi_y=\col X_y$ and $\zeta_y=\col Z_y$.

\begin{lemma}[excess]\label{lem:excess}
For $a\ne b$,
\begin{align*}
  \dim(\xi_a+\zeta_b)+\rk Y_b&\le\dim\zeta_b+r,\\
  \dim(\xi_a+\zeta_b)+\rk W_a&\le\dim\xi_a+s.
\end{align*}
\end{lemma}

\begin{proof}
By Lemma~\ref{lem:slices}, the matrix $M=X_aY_b^{\top}$ is equal to $-Z_bW_a^{\top}$. Its column space lies
in $\xi_a\cap\zeta_b$. By Sylvester's inequality, $\rk M\ge\rk X_a+\rk Y_b-r$ and $\rk M\ge\rk Z_b+\rk W_a-s$.
With $\dim(\xi_a+\zeta_b)=\dim\xi_a+\dim\zeta_b-\dim(\xi_a\cap\zeta_b)$, the first lower bound on $\rk M$ gives
the first inequality, and the second gives the second.
\end{proof}

We use submodularity of dimension in the following form: for subspaces $E_1$, $E_2$, $E_3$ of $\F^n$,
\begin{equation}\label{eq:submod}
  \dim(E_1+E_3)+\dim E_2\le\dim(E_1+E_2)+\dim(E_2+E_3).
\end{equation}
This holds because $E_1+E_3\subseteq(E_1+E_2)+(E_2+E_3)$ and $E_2\subseteq(E_1+E_2)\cap(E_2+E_3)$.

\begin{lemma}[three-cycle]\label{lem:cycle}
Let $a,b,c\in[n]$ be distinct. Then
\begin{align*}
  n+\rk W_c+\rk Y_b+\rk W_a&\le\rk X_a+r+2s,\\
  n+\rk Y_c+\rk W_b+\rk Y_a&\le\rk Z_a+2r+s.
\end{align*}
\end{lemma}

\begin{proof}
By \eqref{eq:span}, $\dim(\zeta_a+\xi_a)=n$. Two applications of \eqref{eq:submod} give
\begin{align*}
  n+\dim\xi_c&\le\dim(\zeta_a+\xi_c)+\dim(\xi_c+\xi_a),\\
  \dim(\xi_c+\xi_a)+\dim\zeta_b&\le\dim(\xi_c+\zeta_b)+\dim(\zeta_b+\xi_a).
\end{align*}
By Lemma~\ref{lem:excess}, applied to the pairs $(c,a)$, $(c,b)$ and $(a,b)$,
\begin{align*}
  \dim(\xi_c+\zeta_a)&\le\dim\xi_c+s-\rk W_c,\\
  \dim(\xi_c+\zeta_b)&\le\dim\zeta_b+r-\rk Y_b,\\
  \dim(\xi_a+\zeta_b)&\le\dim\xi_a+s-\rk W_a.
\end{align*}
The sum of the five inequalities is the first claim, since $\dim\xi_a=\rk X_a$. The second claim follows in
the same way with the roles of $(X,Y,r)$ and $(Z,W,s)$ exchanged.
\end{proof}

\begin{proof}[Proof of Theorem~\ref{thm:43}]
Take three distinct indices $1,2,3$; this needs $n\ge3$. Let $x=\sum_{y=1}^3\rk X_y$, and define $\bar y$,
$z$, $w$ in the same way from the matrices $Y_y$, $Z_y$, $W_y$. Apply Lemma~\ref{lem:cycle} to
$(a,b,c)=(1,2,3)$, $(2,3,1)$ and $(3,1,2)$. In each of the two inequalities, every index occurs once in each
position, so the sums over the three rotations are
\[
  3n+2w+\bar y\le x+3r+6s,\qquad 3n+2\bar y+w\le z+6r+3s.
\]
The transposed system exchanges $X$ with $Y$ and $Z$ with $W$, and gives in the same way
\[
  3n+2z+x\le\bar y+3r+6s,\qquad 3n+2x+z\le w+6r+3s.
\]
These are three rotations with four inequalities each. Their sum is
$12n+2(x+\bar y+z+w)\le18(r+s)$. By the end of Section~\ref{sec:slices}, $x+z\ge3n$ and $\bar y+w\ge3n$.
Hence $24n\le18(r+s)$, that is, $4n\le3(r+s)$.
\end{proof}

\begin{remark}[a shorter proof]\label{rem:short}
Theorem~\ref{thm:43} also has a proof with one triple of indices and no rotations. It was found in the review
of a draft of this note, with AI assistance, and it is not formalized. Write $\eta_y=\col Y_y$ and
$\omega_y=\col W_y$. For $y\ne u$, Lemma~\ref{lem:slices} says that the product of $[X_y\mid Z_u]$ and the
transpose of $[Y_u\mid W_y]$ is zero. Both matrices have $r+s$ columns, so
\begin{equation}\label{eq:pair}
  \dim(\xi_y+\zeta_u)+\dim(\eta_u+\omega_y)\le r+s.
\end{equation}
Let $a,b,c$ be distinct. Two applications of \eqref{eq:submod}, which start from $\dim(\xi_a+\zeta_a)=n$, and
two more for the transposed system, which start from $\dim(\eta_a+\omega_a)=n$, give
\begin{align*}
  n+\dim\zeta_b+\dim\xi_c&\le\dim(\xi_a+\zeta_b)+\dim(\xi_c+\zeta_b)+\dim(\xi_c+\zeta_a),\\
  n+\dim\eta_b+\dim\omega_c&\le\dim(\eta_b+\omega_a)+\dim(\eta_b+\omega_c)+\dim(\eta_a+\omega_c).
\end{align*}
By \eqref{eq:pair} for $(y,u)=(a,b)$, $(c,b)$ and $(c,a)$, the sum of the two right sides is at most $3(r+s)$.
Lemma~\ref{lem:slices} for $y=u$ gives $n\le\rk(X_bY_b^{\top})+\rk(Z_bW_b^{\top})\le\dim\eta_b+\dim\zeta_b$, and
in the same way $n\le\dim\xi_c+\dim\omega_c$. Hence $4n\le3(r+s)$.
\end{remark}

\emph{On the proof of Theorem~\ref{thm:1511}.} Take five distinct indices; this needs $n\ge5$. The proof
is a certificate: a list of $189$ inequalities with integer weights. They relate the widths $r$ and $s$, the
dimensions of sums of the column spaces of the slices at the five indices, and the ranks of the matrices
$X_yY_u^{\top}$. Each inequality is of a kind used above: monotonicity and submodularity of dimension,
\eqref{eq:span}, Sylvester's inequality, and the fact from the proof of Lemma~\ref{lem:excess} that the
column space of $X_yY_u^{\top}$ lies in $\xi_y\cap\zeta_u$, for the slice system and for the transposed
system. The weighted sum of the $189$ inequalities is $15n\le11(r+s)$. A linear-programming search by the
agents found the certificate. We do not print it. The Lean file of Section~\ref{sec:lean} proves each of
the $189$ inequalities for every slice system and checks the weighted sum.

\emph{On the proof of Theorem~\ref{thm:6749}.} The proof is a certificate of the same kind for seven
distinct indices; this needs $n\ge7$. The symmetric group of the seven indices permutes the inequalities, so
the certificate is stated for sums over orbits: it has $77$ inequalities between orbit sums, and their
weighted sum is $67n\le49(r+s)$. A linear-programming search by the agents found it, and the Lean file of
Section~\ref{sec:lean} proves it for every slice system. We do not print it.

\section{Proof of Theorem~\ref{thm:n3}}\label{sec:n3}

The formal proof of Theorem~\ref{thm:n3} has about $4\,900$ lines (Section~\ref{sec:lean}). We outline it
and name the parts of the Lean file. Let $n\ge4$ and $S_n=B+C$ with $r=\rk B$ and $s=\rk\R(C)$. By
Theorem~\ref{thm:43}, $r+s\ge\lceil4n/3\rceil\ge n+2$, so we must exclude $r+s=n+2$. For $n\ge7$
Theorem~\ref{thm:43} does this already, and for $n=6$ Theorem~\ref{thm:1511} does. So the argument below
is needed only for $n=4$ and $n=5$. The formal proof does not use Theorem~\ref{thm:1511}, and we outline it
as it is. Since the realignment
exchanges $r$ and $s$, we may assume $s\le r$. The case $s=0$ is impossible because $\rk S_n=n^2$.

\begin{lemma}[one term]\label{lem:one}
If $s\le1$, then $r\ge n(n-1)/2$.
\end{lemma}

\begin{proof}[Outline]
If $\rk\R(C)\le1$, then $C=G\otimes H$ for two $n\times n$ matrices. The kernel of $S_n-G\otimes H$ is
isomorphic to the space of $n\times n$ matrices $Z$ with $Z^{\top}=GZH^{\top}$, and this space has dimension at
most $n(n+1)/2$ over every field. Hence $r=\rk(S_n-G\otimes H)\ge n^2-n(n+1)/2$.
\end{proof}

Since $n(n-1)/2>n+1$ for $n\ge4$, Lemma~\ref{lem:one} excludes $s=1$.

\begin{lemma}[two terms]\label{lem:two}
If $s\le2$, then $r\ge n+1$.
\end{lemma}

\begin{proof}[Outline]
Pad the slices with zero columns to widths $(n,2)$. One shows that every $W_y$ is injective and that every
$X_i$ has a nonzero vector $w_i$ with $w_iX_i=0$. Lemma~\ref{lem:slices} then gives $w_iZ_u=0$ for $u\ne i$
and $w_iZ_i\ne0$. For a fixed $u$, the map $g\mapsto\sum_ig_iw_i$ is injective on the vectors $g$ with
$g_u=0$, and its image lies in the left kernel of $Z_u$. So $\rk Z_u\le1$, which contradicts the injectivity
statements for the transposed system.
\end{proof}

Lemma~\ref{lem:two} excludes $s=2$, because then $r=n<n+1$. So $3\le s\le r$, and $r=n+2-s\le n-1$.

\begin{lemma}[counting]\label{lem:count}
Let $n\ge5$, $r+s=n+2$ and $3\le s\le r$. There is no slice system of widths $(r,s)$.
\end{lemma}

\begin{proof}[Outline]
First, every $W_z$ and, by the transposed system, every $Z_u$ has full rank $s$. If $\rk X_z=r$, this
follows from a dimension count with the subspace $\ker W_z^{\top}$. If $\rk X_z=r-1$, the projection
$X_zY_z^{\top}$ onto $\col X_z$ turns the system into a system on the other $n-1$ indices with a term of rank
one, and a count of ranks gives $(n-1)(r-1)\le(r-1)^2+s$, which is false. At $n=5$ the shape is $(4,3)$, and
the same conclusion has a separate proof. Second, every $X_i$ is an $n\times r$ matrix with $r<n$, so it has
a nonzero left kernel vector. As in the proof of Lemma~\ref{lem:two}, these two facts contradict each other.
\end{proof}

\begin{lemma}[$n=4$]\label{lem:four}
There is no slice system with $n=4$ and widths $(3,3)$.
\end{lemma}

\begin{proof}[Outline]
Let $K=\F(t_1,\dots,t_4)$ and consider the $4\times3$ matrices $X(t)=\sum_at_aX_a$ and $W(t)=\sum_at_aW_a$
over $K$. Two facts about the generic point $t$ are used: it satisfies no nonzero $\F$-linear relation, and
an $\F$-matrix that has $t$ as an eigenvector is scalar. If $\rk X(t)=3$, a comparison of images, together
with the fact that the matrices $Y_u$ have no common nonzero kernel vector, gives a contradiction; the same
holds for the transposed system. Otherwise both ranks are at most $2$, and then also at every point of
$\F^4$. Lemma~\ref{lem:slices} gives $W(x)^{\top}X(x)=0$ and $X(x)\,Y(y)^{\top}X(x)=(x\cdot y)\,X(x)$ for all
$x,y\in\F^4$, where $Y(y)=\sum_by_bY_b$, and these identities force a contradiction through the kernel
vectors of $X(x)$.
\end{proof}

\begin{proof}[Proof of Theorem~\ref{thm:n3}]
By the discussion above, $r+s=n+2$ forces $3\le s\le r$. Lemma~\ref{lem:four} excludes this for $n=4$, and
Lemma~\ref{lem:count} for $n\ge5$.
\end{proof}

\section{Small cases and remarks}\label{sec:small}

\begin{proof}[Proof of Corollary~\ref{cor:values}]
The lower bounds follow from Theorems~\ref{thm:43}, \ref{thm:1511}, \ref{thm:6749} and~\ref{thm:n3}, since
$67n/49\ge15n/11$ for all $n$ and $67n/49>n+2$ for $n\ge6$; at $n=3$, Theorem~\ref{thm:43} gives
$r+s\ge4$. For the upper bounds, split $S_n=(S_n-I)+I$. The realignment of $I$ is
$\R(I)[(x,u),(y,v)]=[x=u]\,[y=v]$, a matrix of rank one. The matrix $S_n$ is the permutation matrix of the
involution $(x,y)\mapsto(y,x)$, which has $n$ fixed points and $n(n-1)/2$ orbits of size two. So $S_n-I$ is,
after a permutation of rows and columns, block diagonal with $n$ zero blocks of size one and $n(n-1)/2$ blocks
$\left(\begin{smallmatrix}-1&1\\1&-1\end{smallmatrix}\right)$ of rank one, over every field. Hence
$P_\F(n)\le n(n-1)/2+1$, which is $4$ at $n=3$, $7$ at $n=4$ and $11$ at $n=5$.
\end{proof}

\begin{remark}[small $n$]
The same splitting gives $P_\F(2)\le2$, so $P_\F(2)=2$. Trivially $P_\F(1)=1$.
\end{remark}

\begin{remark}[a border splitting at $n=4$]\label{rem:border}
The multi-agent run described below also produced $16\times3$ matrices $U(t),V(t),G(t),H(t)$ whose entries
are integer polynomials of degree at most $3$ in $t$, with the following property. Let $B(t)=U(t)V(t)^{\top}$,
with rows indexed by $(x,y)$ and columns by $(u,v)$, and let $C(t)$ be the matrix with
$\R(C(t))=G(t)H(t)^{\top}$. Then
\[
  B(t)+C(t)=c\,t^3S_4+O(t^4),
\]
with $c=1994028894245376=2^9\cdot3^7\cdot29\cdot59\cdot71\cdot107\cdot137$.
Let $\F$ be a field in which $c\ne0$, that is, of characteristic $0$ or of a characteristic other than $2$, $3$, $29$, $59$, $71$, $107$
and $137$. Then $S_4$ lies in the Zariski closure of the set of sums of a matrix of rank at
most $3$ and a matrix whose realignment has rank at most $3$; over $\mathbb R$ or $\mathbb C$ it is the limit
of $(B(t)+C(t))/(ct^3)$ as $t\to0$. So the border analogue of $P_\F(4)$ is at most $6$, while
$P_\F(4)=7$ by Corollary~\ref{cor:values}: the proof of Lemma~\ref{lem:four} cannot pass to limits. The
trivial bound survives limits, because rank is lower semicontinuous, so the border value at $n=4$ lies
between $4$ and $6$. We do not know whether Theorem~\ref{thm:43} survives limits; if it does, the border
value at $n=4$ is $6$. The data and a script that checks the identity with exact integer arithmetic are in
the repository of Section~\ref{sec:lean}, in the files \texttt{border\_n4.json} and
\texttt{verify\_border\_n4.py}.
\end{remark}

\begin{remark}[width versus bilinear complexity]
For $n\in\{4,8\}$ the four-square and eight-square identities have integer coefficients, so $\beta(n)\le n$,
while $P_\Q(4)=7$ and $P_\Q(8)\ge11$ by Corollary~\ref{cor:values}. Hence the lower inequality
$\beta(n)\le P_\Q(n)$ of \cite[Proposition~4.1]{HWY11} is strict for $n=4$ and $n=8$.
\end{remark}

\begin{remark}[limits of the arguments]\label{rem:limits}
The proofs of Theorems~\ref{thm:43}, \ref{thm:1511} and~\ref{thm:6749} use only facts of the following
kinds, for three, five or seven indices. For the column spaces $\xi_y$, $\zeta_y$, $\eta_y=\col Y_y$ and
$\omega_y=\col W_y$: \eqref{eq:span}, Lemma~\ref{lem:excess}, \eqref{eq:pair}, and inequalities between
dimensions that hold for all subspaces. For the ranks of the matrices $X_yY_u^{\top}$ with $y\ne u$: the two
lower bounds from Sylvester's inequality and the inclusion of the column space in $\xi_y\cap\zeta_u$, as in
the proof of Lemma~\ref{lem:excess}. And the same facts for the transposed system. With $L$ indices,
computer searches for the best combination of such inequalities give the constant $4/3$ for $L=3$ and
$L=4$, $15/11$ for $L=5$ (Theorem~\ref{thm:1511}) and $67/49\approx1.367$ for $L=7$
(Theorem~\ref{thm:6749}). A search by the agents and an independent search in the review of a draft agree
on the values for $L\le5$. Further computations by the agents suggest that the constants increase with $L$
to the limit $1+1/e\approx1.368$; we have not checked this.

No such combination reaches the constant $2$, by the following example from the same review. Let $L\ge3$
divide $n$, split $\F^n$ into $L$ blocks of dimension $n/L$, and put $r=(2-3/L)\,n$ and $s=n/L$. For
$y\in[L]$ let $\xi_y=\eta_y$ be the sum of all blocks except block $y$, and let $\zeta_y=\omega_y$ be block
$y$. Then $\xi_y+\zeta_y=\F^n$, all dimensions are at most the widths, and for $a\ne b$ we have
$\zeta_b\subseteq\xi_a$, so the two inequalities of Lemma~\ref{lem:excess} and \eqref{eq:pair} hold, with
equality. For $a\ne b$ give $X_aY_b^{\top}$ the rank $n/L$ and the column space $\zeta_b$, which is
$\xi_a\cap\zeta_b$; then both Sylvester bounds, $\rk X_a+\rk Y_b-r$ and $\rk Z_b+\rk W_a-s$, are equal to
$n/L$, so they hold with equality, and so does the inclusion. The same is true for the transposed system by
symmetry. We do not claim that these subspaces and ranks come from a splitting, but every argument that
uses only the facts above for $L$ indices also applies to them. So no such argument proves more than
$r+s\ge(2-2/L)\,n$, and none gives a superlinear bound. The proof of Theorem~\ref{thm:n3} excludes
the single total $n+2$ and does not extend directly. The best upper bounds we know are $1+n(n-1)/2$ from the
splitting above, $P_\Q(8)\le16$ from the eight-square identity~\cite[Appendix~G]{Cla26}, and
$O(n^{1.62})$~\cite[Theorem~12]{Hru24}.
\end{remark}

\section{Formal verification}\label{sec:lean}

{\sloppy The repository \url{https://github.com/Arth-Singh/crossing-lower-bounds} contains four Lean files
with the proofs. Everything in this section refers to the tag \texttt{lean-v1} of the repository (commit
\texttt{3d302668a9f3}); the Comparator check of the Palomar registry passed on this commit. They use Lean~4~\cite{Lean4} and Mathlib~\cite{Mathlib} only, and the repository pins Lean
and Mathlib v4.35.0-rc2. The file \texttt{SwapSplit.lean} ($304$ lines) defines the swap matrix and the
realignment as above and proves Theorem~\ref{thm:43} as \texttt{SwapSplit.swap\_split\_ge\_four\_thirds}, with
the proof of Section~\ref{sec:43}. The file \texttt{SmallValues.lean} ($4\,922$ lines) proves
Theorem~\ref{thm:n3} as \texttt{SwapSplit.swap\_split\_ge\_add\_three}, the value $P_\F(4)=7$, including the
upper bound, as \texttt{SwapSplit.swap\_split\_four\_eq\_seven}, and $P_\F(5)\ge8$ and $P_\F(6)\ge9$. Its body
is the dependency closure of $61$ files of the multi-agent run, unchanged, and its last section proves that
the agents' definitions of the swap matrix and the realignment are those of \texttt{SwapSplit.lean}. The file
\texttt{FifteenElevenths.lean} ($1\,761$ lines, with $17$ more files of the run) proves
Theorem~\ref{thm:1511} as \texttt{SwapSplit.swap\_split\_ge\_fifteen\_elevenths}, and the file
\texttt{SixtySevenFortyNinths.lean} ($1\,294$ lines, with $11$ more files) proves Theorem~\ref{thm:6749}
as \texttt{SwapSplit.swap\_split\_ge\_sixtyseven\_fortyninths}. The files \texttt{Challenge.lean} and
\texttt{Solution.lean} restate the seven results for a Comparator check: \texttt{Challenge.lean} has
$88$ lines, imports one Mathlib module, and contains the two definitions and the seven statements. Every
theorem named here depends only on the axioms \texttt{propext}, \texttt{Classical.choice} and
\texttt{Quot.sound}. The upper bounds at $n=3$ and $n=5$ are not formalized, and the identity of
Remark~\ref{rem:border} is checked by the script \texttt{verify\_border\_n4.py}.\par}

The definitions and the statements of Theorems~\ref{thm:43}, \ref{thm:1511}, \ref{thm:6749}
and~\ref{thm:n3} and of $P_\F(4)=7$ are as follows, verbatim from the Lean files, without the proofs. Here \texttt{Matrix m n F} is the type of functions
\texttt{m → n → F}, and \texttt{rank} is Mathlib's \texttt{Matrix.rank}.

{\footnotesize
\begin{verbatim}
def swapMat (F : Type*) [Field F] (N : ℕ) : Matrix (Fin N × Fin N) (Fin N × Fin N) F :=
  fun p q => if q.1 = p.2 ∧ q.2 = p.1 then 1 else 0

def realign {F : Type*} [Field F] {N : ℕ} (C : Matrix (Fin N × Fin N) (Fin N × Fin N) F) :
    Matrix (Fin N × Fin N) (Fin N × Fin N) F :=
  fun p q => C (p.1, q.1) (p.2, q.2)

theorem swap_split_ge_four_thirds {F : Type*} [Field F] (N : ℕ) (hN : 3 ≤ N)
    (B C : Matrix (Fin N × Fin N) (Fin N × Fin N) F)
    (h : swapMat F N = B + C) : 4 * N ≤ 3 * (B.rank + (realign C).rank)

theorem swap_split_ge_fifteen_elevenths {F : Type*} [Field F] (N : ℕ) (hN : 5 ≤ N)
    (B C : Matrix (Fin N × Fin N) (Fin N × Fin N) F)
    (h : swapMat F N = B + C) : 15 * N ≤ 11 * (B.rank + (realign C).rank)

theorem swap_split_ge_sixtyseven_fortyninths {F : Type*} [Field F] (N : ℕ) (hN : 7 ≤ N)
    (B C : Matrix (Fin N × Fin N) (Fin N × Fin N) F)
    (h : swapMat F N = B + C) : 67 * N ≤ 49 * (B.rank + (realign C).rank)

theorem swap_split_ge_add_three {F : Type*} [Field F] (N : ℕ) (hN : 4 ≤ N)
    (B C : Matrix (Fin N × Fin N) (Fin N × Fin N) F)
    (h : swapMat F N = B + C) : N + 3 ≤ B.rank + (realign C).rank

theorem swap_split_four_eq_seven (F : Type*) [Field F] :
    (∃ B C : Matrix (Fin 4 × Fin 4) (Fin 4 × Fin 4) F,
        swapMat F 4 = B + C ∧ B.rank + (realign C).rank = 7) ∧
    ∀ B C : Matrix (Fin 4 × Fin 4) (Fin 4 × Fin 4) F,
        swapMat F 4 = B + C → 7 ≤ B.rank + (realign C).rank
\end{verbatim}
}

\section{How the results were found}\label{sec:how}

The results come from one run of an in-house harness in which about $20$ AI agents (Claude Opus~5.5 workers and a
Claude Fable~5.2 lead, Anthropic) work as a small research organization. Agents hold roles (leads, workers, an advisor), set goals
for each other, and share a library of Lean files. A statement counts as proved only after Lean accepts it,
and a claim of novelty needs a literature search by the advisor. The run studied the decomposition of the
crossing that is used in~\cite{Cla26}.

The steps were as follows. One worker agent posed the lower-bound question as a goal and proved the trivial
bound $n$ in Lean. A second worker proved $r+s\ge n+1$ for $n\ge3$, and then $7(r+s)\ge8n$, with the
identities of Lemma~\ref{lem:slices}; no person had asked for these bounds. The author then gave a small team
one hour to go beyond them, and the first worker proved $r+s\ge n+2$ for $n\ge4$ and $4(r+s)\ge5n$ within
$36$ minutes. The author then put the whole organization on the problem for $35$ minutes. After $15$
minutes, a third worker proved Theorem~\ref{thm:43}, with the excess lemma in rank form and the transposed
system; another worker found the same constant independently, as a combination of $26$ inequalities found by
a linear-programming search. The value $P_\F(5)\ge8$ followed after $19$ minutes, $P_\Q(4)=7$ after $24$
minutes, $P_\F(4)=7$ over every field after $33$ minutes, and $P_\F(6)\ge9$, which completes
Theorem~\ref{thm:n3}, after $44$ minutes, in an extension of the time. In a second phase, aimed at
superlinear bounds, the agents extended the slice argument to five and seven indices.
Theorem~\ref{thm:1511} followed after $62$ minutes and Theorem~\ref{thm:6749} after $74$ minutes. For each, one
worker found the constant by a linear-programming search, other agents made the certificate exact, and
others proved its inequalities in Lean. The earlier bounds are consequences of
Theorem~\ref{thm:43} and are not included here. Another worker found the border splitting of
Remark~\ref{rem:border}. The advisor agent refereed the proofs, checked the border splitting with exact
arithmetic, and searched the literature for earlier bounds.

The Lean files of the repository, the literature check for this note and the text of this note were
prepared with AI assistance (Claude). Theorems~\ref{thm:43}, \ref{thm:1511}, \ref{thm:6749} and~\ref{thm:n3}
are verified in Lean (Section~\ref{sec:lean}). The author takes responsibility for the content.

\begin{thebibliography}{Mathlib}

\bibitem[BKW26]{BKW26}
C.~Brand, P.~Kaski and J.~Wang, \emph{Partition rank and algebraic circuit lower bounds}, arXiv:2607.02241,
2026.

\bibitem[Cla26]{Cla26}
Claude (Anthropic), \emph{A $2^{o(n)}$ time algorithm for 3-SAT}, Hexagon preprint hexagon:2610.00577v1,
October 2026. Lean files: \url{https://github.com/anthropics/formal-math} (folder \texttt{eth}).

\bibitem[Hru24]{Hru24}
P.~Hrube\v{s}, \emph{A subquadratic upper bound on sum-of-squares composition formulas}, in: 39th
Computational Complexity Conference (CCC 2024), LIPIcs 300, 12:1--12:11, 2024,
doi:10.4230/LIPIcs.CCC.2024.12.

\bibitem[HWY11]{HWY11}
P.~Hrube\v{s}, A.~Wigderson and A.~Yehudayoff, \emph{Non-commutative circuits and the sum-of-squares
problem}, J. Amer. Math. Soc. \textbf{24} (2011), no.~3, 871--898, doi:10.1090/S0894-0347-2011-00694-2.

\bibitem[Kar23]{Kar23}
T.~Karam, \emph{The interplay between bounded ranks of tensors arising from partitions}, arXiv:2310.11356,
2023.

\bibitem[Lean4]{Lean4}
L.~de~Moura and S.~Ullrich, \emph{The Lean~4 theorem prover and programming language}, in: Automated
Deduction -- CADE 28, LNCS 12699, Springer, 2021, 625--635.

\bibitem[Mathlib]{Mathlib}
The mathlib Community, \emph{The Lean mathematical library}, in: Proceedings of the 9th ACM SIGPLAN
International Conference on Certified Programs and Proofs (CPP 2020), 367--381.

\bibitem[Nas20]{Nas20}
E.~Naslund, \emph{The partition rank of a tensor and $k$-right corners in $\mathbb{F}_q^n$}, J. Combin.
Theory Ser. A \textbf{174} (2020), 105190, arXiv:1701.04475.

\end{thebibliography}

\end{document}
